\chapter{Revue de litterature}

\begin{simplebox2}
	Dans cette partie, nous présentons la revue de littérature relative au suivi et à l’évaluation des Objectifs de Développement Durable (ODD).  
	Elle comprend trois volets : la revue théorique, qui expose les concepts fondamentaux du développement durable ; la revue empirique, qui analyse les principales études et résultats observés ; et la revue méthodologique, qui s’intéresse aux approches et outils mobilisés pour la mesure des indicateurs.  
	L’objectif est de dégager les apports et limites des travaux existants afin d’orienter la démarche méthodologique adoptée dans le cadre de cette étude.
\end{simplebox2}

	

\section{Revue théorique}
Depuis l’an 2000, la communauté internationale s’est dotée d’un cadre pour le développement à court terme, les \textit{Objectifs du Millénaire pour le Développement} (OMD), adoptés lors du Sommet du Millénaire des Nations Unies \autocite{un2000}. Ces huit objectifs portaient sur la réduction de l’extrême pauvreté et de la faim, l’éducation primaire universelle, l’égalité des sexes, la réduction de la mortalité infantile, l’amélioration de la santé maternelle, la lutte contre le VIH/Sida, le paludisme et d’autres maladies, la préservation de l’environnement et enfin la mise en place d’un partenariat mondial pour le développement \autocite{who2015}. Ces OMD constituaient une avancée majeure car ils proposaient des cibles chiffrées à atteindre d’ici 2015 et servaient de cadre de suivi pour les politiques nationales et internationales.

Au Sénégal, des progrès significatifs ont été enregistrés dans la mise en œuvre des OMD. Concernant la lutte contre la faim, le pays a réussi à ramener le taux de malnutrition de près de 24,5~\% à environ 10~\% entre 1992 et 2015, atteignant ainsi la cible de réduction de moitié \autocite{fao2015}. Dans le domaine de l’accès à l’eau potable, le Sénégal a également progressé : le taux de desserte en milieu urbain a dépassé la cible OMD et l’accès national à l’eau atteignait environ 88,5~\% en 2015 \autocite{bm2015}. Toutefois, certaines limites sont restées importantes, notamment dans la réduction de la pauvreté monétaire et du chômage, ainsi que dans l’atteinte de l’éducation primaire universelle et de la réduction de la mortalité maternelle \autocite{lo2015}.

L’expérience des OMD a mis en évidence plusieurs contraintes structurelles : des financements insuffisants, des disparités régionales, une faible capacité logistique et des déficits dans la collecte des données statistiques, qui ont freiné l’atteinte de certaines cibles \autocite{ipar2015}. De plus, certains indicateurs n’étaient pas toujours désagrégés par genre ou par zone géographique, rendant l’évaluation incomplète \autocite{eca2015}. Ces limites ont nourri la réflexion ayant conduit à l’adoption en 2015 des \textit{Objectifs de Développement Durable} (ODD), plus larges et inclusifs, fixés à l’horizon 2030.

Les enseignements tirés des OMD montrent qu’il est essentiel de disposer d’indicateurs nationaux clairs, adaptés aux réalités locales et suffisamment robustes pour mesurer les progrès. Ainsi, pour le Sénégal, la construction d’un indicateur national ODD s’appuie sur ce bilan contrasté : succès en matière d’accès à l’eau et de lutte contre la faim, mais difficultés persistantes en matière de pauvreté monétaire, d’éducation et de santé. En ce sens, les OMD constituent une base théorique incontournable pour comprendre les défis actuels liés au suivi des ODD et justifient l’importance de mettre en place des indicateurs pertinents, fiables et contextualisés.

%=================================================================
\section{Revue empirique}

Les résultats empiriques disponibles montrent des progrès notables du Sénégal dans plusieurs domaines des ODD, mais également des défis persistants.

Sur le plan social, les enquêtes \textit{EDS-C 2023} confirment la baisse continue de la mortalité infantile et juvénile : le taux de mortalité néonatale est estimé à 22 décès pour 1~000 naissances vivantes, tandis que la mortalité des moins de 5 ans se situe à 44 pour 1~000 \autocite{ANSD_EDSC_2023}. En matière de santé maternelle, le ratio de mortalité maternelle est de 236 décès pour 100~000 naissances vivantes. La couverture vaccinale reste élevée : 85~\% des enfants de 12–23 mois ont reçu les vaccins de base \autocite{ANSD_EDSC_2023}.

Du côté de l’éducation, le \textit{Sustainable Development Report 2025} indique un taux net de scolarisation primaire de 75,9~\% en 2023 et un taux d’achèvement du secondaire inférieur de 39,3~\% \autocite{SDR2025}, traduisant à la fois des progrès et des effets liés au redoublement et à la correction des effectifs. Le taux d’alphabétisation des jeunes (15–24 ans) atteignait 78,1~\% en 2020 \autocite{ANSD_EDSC_2023}.

Concernant l’égalité de genre, l’\textit{ENR-VFFS 2023} révèle qu’environ 7~\% des femmes de 15–49 ans déclarent avoir subi une forme de violence sexuelle au cours de leur vie, et qu’une femme sur dix a été mariée avant l’âge de 15 ans \autocite{ANSD_ENRVFFS_2023}. Du point de vue de la représentation politique, les femmes n’occupent que 19,9~\% des sièges au Parlement \autocite{SDR2025}.

Les indicateurs économiques témoignent d’avancées contrastées. Le Sénégal affiche une croissance soutenue avec un PIB par habitant en progression de 3,7~\% en moyenne annuelle sur la période 2015–2022 \autocite{BanqueMondiale2022}. Toutefois, le chômage reste préoccupant, touchant 21,3~\% des jeunes, et le taux de pauvreté monétaire demeure autour de 37,8~\% \autocite{ANSD_EHCVM_2018_2019}. Les envois de fonds des travailleurs migrants représentent un levier majeur : ils équivalaient à 9,6~\% du PIB en 2022 \autocite{BM_KNOMAD_2022}.

Sur le plan environnemental, les données du SDR 2025 indiquent des défis persistants. Les émissions de CO\textsubscript{2} liées à la combustion des énergies fossiles atteignaient 0,6~t par habitant en 2021 \autocite{SDR2025}. La proportion de la population utilisant au moins des services de base d’eau potable est estimée à 98,7~\% en 2022, contre 95,3~\% pour l’assainissement de base. En revanche, le prélèvement d’eau douce représente 97,4~\% des ressources disponibles \autocite{SDR2025}.

En matière de gouvernance et d’institutions, le Sénégal affiche un taux d’enregistrement des naissances de 92~\% des enfants de moins de 5 ans, mais des marges de progrès existent pour la lutte contre la corruption (indice de perception à 43/100 en 2024) \autocite{Transparency2024}.

Enfin, sur le plan de la coopération internationale, les dépenses publiques combinées en santé et éducation représentaient 8,7~\% du PIB en 2023, tandis que les financements concessionnels reçus restaient en deçà des cibles internationales (0,5~\% du RNB en 2023) \autocite{SDR2025}.

\bigskip
Ces données empiriques fournissent une base solide pour la construction d’un indicateur ODD national au Sénégal, intégrant des composantes sociales, économiques, environnementales et institutionnelles, tout en soulignant les lacunes à combler pour obtenir une mesure complète et actualisée.









\section{Revue méthodologique}


Les indicateurs ODD sont choisis non seulement pour leur alignement avec le cadre global des Nations Unies, mais également pour leur capacité à refléter les enjeux socio-économiques propres au contexte sénégalais et à être mesurés de manière régulière à partir de sources statistiques fiables (enquêtes, recensements, registres administratifs, rapports sectoriels).

\subsection*{Choix des indicateurs}

Les indicateurs retenus sont choisis selon trois critères principaux : 
\begin{simplebox} 
\begin{enumerate}
\item \textbf{Pertinence} : l’indicateur doit refléter fidèlement le phénomène qu’il est censé mesurer.  

\item \textbf{Disponibilité et fiabilité} : les données doivent être accessibles régulièrement et issues de sources statistiques crédibles.  

\item \textbf{Comparabilité} : les indicateurs doivent permettre des comparaisons dans le temps et entre pays ou régions.  

Cette sélection permet de disposer d’un suivi robuste et cohérent des progrès vers l’Agenda 2030, tout en restant adapté aux priorités nationales.
\end{enumerate}
\end{simplebox}


\subsection*{Pondération des indicateurs}
Un enjeu méthodologique central réside dans la question des pondérations à appliquer lors de la construction d’un indice ODD composite. La manière dont les indicateurs sont pondérés influence directement l’importance relative accordée à chaque dimension et, par conséquent, les résultats obtenus. La littérature distingue généralement quatre grandes approches. La première est celle de la pondération égale, qui consiste à attribuer le même poids à tous les objectifs et à tous les indicateurs. Cette méthode, retenue par le \textit{SDG Index and Dashboards}, repose sur l’idée que chaque objectif revêt une importance équivalente dans le cadre de l’Agenda 2030. Son principal avantage est sa simplicité et sa transparence, mais elle introduit une pondération implicite : les objectifs comprenant de nombreux indicateurs (comme la santé ou l’éducation) voient chaque mesure individuelle peser relativement moins que dans les objectifs ne reposant que sur un nombre réduit d’indicateurs.

La deuxième approche est celle des pondérations dites « mathématiques », souvent dérivées de méthodes statistiques comme l’Analyse en Composantes Principales (ACP) ou les charges factorielles. Elles permettent d’attribuer des poids en fonction de la corrélation entre variables, en privilégiant celles qui contribuent le plus à un facteur sous-jacent commun. Bien qu’intéressante sur le plan analytique, cette méthode se heurte à l’hétérogénéité intrinsèque des ODD. Par exemple, l’ODD~2 sur la sécurité alimentaire regroupe des dimensions très diverses – nutrition, obésité, agriculture durable – qui ne présentent pas nécessairement de corrélation statistique. Dans de tels cas, l’application d’une pondération mathématique peut conduire à des résultats peu interprétables au niveau national.

Une troisième option repose sur les pondérations d’experts. Celles-ci consistent à mobiliser des spécialistes pour déterminer, par consensus, les dimensions jugées prioritaires et leur attribuer des poids plus ou moins élevés. Dans le contexte sénégalais, une telle démarche pourrait refléter les priorités nationales, par exemple en accordant une importance particulière à la réduction de la pauvreté, à l’éducation ou à l’égalité de genre. Toutefois, cette méthode reste subjective et dépend de l’adhésion des communautés d’experts, qui ne parviennent pas toujours à un accord.

Enfin, certains indices internationaux expérimentent des pondérations subjectives ou flexibles, laissant aux utilisateurs la possibilité de définir eux-mêmes les poids à accorder à chaque dimension. Cette méthode a été popularisée par l’Indice du bien-être de l’OCDE, mais elle paraît peu adaptée aux ODD. En effet, elle risque d’encourager une lecture sélective, où les objectifs jugés difficiles à atteindre seraient minimisés au profit de ceux qui sont plus accessibles. Dans le cadre d’un suivi institutionnel, cette approche n’offre donc pas les garanties nécessaires de neutralité et de comparabilité.

En définitive, la pondération égale demeure la solution la plus répandue et la plus adaptée pour le suivi des ODD au Sénégal. Elle assure une neutralité méthodologique et une meilleure lisibilité des résultats, tout en restant conforme aux pratiques internationales. Il convient néanmoins de souligner que l’égalité des poids ne signifie pas l’absence de pondération : la différence du nombre d’indicateurs entre objectifs introduit une pondération implicite qui peut affecter l’équilibre global de l’indice. C’est pourquoi le renforcement progressif de la disponibilité des données constitue un enjeu crucial afin de réduire ces déséquilibres et d’améliorer la robustesse des indicateurs produits.




















\subsection*{Agrégation des indicateurs}

Après la pondération, se pose la question de l’agrégation, c’est-à-dire la manière de combiner les différents indicateurs pour obtenir un indice global. Ce choix méthodologique est déterminant car il conditionne la lecture finale des résultats. La littérature distingue principalement trois approches.

\paragraph{1. Moyenne arithmétique (durabilité faible)}  
Cette approche suppose que les indicateurs sont parfaitement substituables. Une faiblesse sur un indicateur peut être compensée par une bonne performance sur un autre. Par exemple, un retard dans l’éducation pourrait être « neutralisé » par de bons résultats en santé. Cette méthode est simple, lisible et largement utilisée. La formule correspondante est :
\begin{simplebox2}
\begin{equation}
I_j = \frac{1}{N_j} \sum_{k=1}^{N_j} I_{jk}
\end{equation}
où $I_j$ est l’indice du SDG $j$, $N_j$ le nombre d’indicateurs pour cet objectif, et $I_{jk}$ le score de l’indicateur $k$.
\end{simplebox2}

\paragraph{2. Minimum de Leontief (durabilité forte)}  
Ici, aucun indicateur ne peut compenser un autre. L’indice global est déterminé par la plus faible valeur observée parmi les indicateurs :
\begin{simplebox2}
\begin{equation}
I_j = \min \{ I_{jk} \}_{k=1,\dots,N_j}
\end{equation}
Cette méthode reflète mieux l’esprit des ODD, considérés comme indivisibles, mais elle peut masquer les progrès réalisés dans d’autres domaines.
\end{simplebox2}

\paragraph{3. Moyenne géométrique (substituabilité intermédiaire)}  
Cette approche est intermédiaire et limite l’effet de compensation excessive entre indicateurs. Elle est utilisée, par exemple, pour l’Indice de Développement Humain (IDH) :
\begin{simplebox2}
\begin{equation}
I_j = \left( \prod_{k=1}^{N_j} I_{jk} \right)^{1/N_j}
\end{equation}
Une très faible performance sur un indicateur tire mécaniquement l’ensemble vers le bas. Elle est adaptée lorsque les indicateurs sont hétérogènes et que l’on souhaite éviter la neutralisation des faibles scores par des scores élevés.
\end{simplebox2}

\paragraph{Choix pour le présent rapport}  
Pour le suivi des ODD au Sénégal, l’agrégation repose sur un compromis entre simplicité, rigueur et lisibilité. La moyenne arithmétique a été retenue pour calculer les indices ODD, car elle permet une lecture claire des résultats, une comparabilité internationale, tout en restant facile à mettre en œuvre avec des données nationales. Chaque indicateur est pondéré de manière égale, et le poids relatif de chaque indicateur dans un objectif est inversement proportionnel au nombre d’indicateurs considérés pour cet objectif.





\subsection*{Note sur l’agrégation avancée (CES)}  
Pour information méthodologique, il existe également des fonctions d’agrégation plus sophistiquées, telles que la fonction \textit{constant-elasticity-of-substitution} (CES), qui permettent de moduler le degré de substituabilité entre indicateurs :
\begin{equation}
I_j = \left[ \sum_{k=1}^{N_j} I_{jk}^{-\rho} \right]^{-1/\rho}, \quad -1 \leq \rho \leq \infty
\end{equation}
où $\rho$ contrôle la substituabilité : $\rho=-1$ correspond à une substituabilité parfaite (moyenne arithmétique), $\rho=\infty$ à une substituabilité nulle (Leontief). La moyenne géométrique correspond à une valeur intermédiaire de $\rho$. Cette approche, bien que plus flexible, n’a pas été utilisée dans ce rapport afin de privilégier la lisibilité et la simplicité d’interprétation.
